% Source: physical PDF pp. 155--167 (printed pp. 132--144), from the Section
% 27.4 heading through the opening paragraph on physical p. 167 immediately
% before Section 27.5.
% Inventory: Eqs. (27.4.1)--(27.4.42), seven unnumbered display groups, one
% linked citation, two ordinary footnotes, one centered three-asterisk divider,
% and no figures or tables.
% Source note: the intermediate unnumbered spinor identity and the extracted
% F-term visibly print \psi_B where the surrounding formulas use \lambda_B.
% The theta_L^2 expansion and the same identity visibly omit D_A in the
% relevant term, while the identity visibly omits f_{A\mu\nu} on its right.
% The unnumbered M_0^2 relation after (27.4.19) also visibly repeats t_B in
% the vector on the right despite its sum over A. These printed forms are
% preserved.
% Uncertainties: none.

\subsection{Renormalizable Gauge Theories with Chiral Superfields}

We will now put together the pieces assembled in the previous three sections,
to construct the most general renormalizable action for chiral superfields
interacting with general gauge fields. Adding the terms
\eqref{eq:27.1.27}, \eqref{eq:27.2.7}, and \eqref{eq:27.3.1} and the
superpotential terms in Eq.~\eqref{eq:26.4.5} gives the Lagrangian density
\begin{equation}
  \begin{aligned}
  \mathcal{L}={}&
  \frac12\left[
    \Phi^\dagger
    \exp\left(-2\sum_A t_A V_A\right)\Phi
  \right]_D
  -\frac12\operatorname{Re}\sum_A
    \left(W_{AL}^{\mathrm T}\epsilon W_{AL}\right)_{\mathcal{F}}
  \\
  &-\frac{g^2\theta}{16\pi^2}\sum_A
    \operatorname{Im}
    \left(W_{AL}^{\mathrm T}\epsilon W_{AL}\right)_{\mathcal{F}}
  -\sum_A\xi_A[V_A]_D
  +2\operatorname{Re}[f]_{\mathcal{F}}
  \\
  ={}&
  -\sum_n(D_\mu\phi)_n^*(D^\mu\phi)_n
  -\frac12\sum_n
    \left(\bar\psi_n\gamma^\mu(D_\mu\psi_L)_n\right)
  +\sum_n\mathcal{F}_n^*\mathcal{F}_n
  \\
  &-\operatorname{Re}\sum_{nm}
    \frac{\partial^2f(\phi)}{\partial\phi_n\partial\phi_m}
    \left(\psi_{nL}^{\mathrm T}\epsilon\psi_{mL}\right)
  +2\operatorname{Re}\sum_n
    \frac{\partial f(\phi)}{\partial\phi_n}\mathcal{F}_n
  \\
  &-2\sqrt{2}\operatorname{Im}\sum_{Anm}
    (t_A)_{nm}\left(\bar\psi_{nL}\lambda_A\right)\phi_m
  +2\sqrt{2}\operatorname{Im}\sum_{Anm}
    (t_A)_{mn}\left(\bar\psi_{nR}\lambda_A\right)\phi_m^*
  \\
  &-\sum_{Anm}\phi_n^*(t_A)_{nm}\phi_mD_A
  -\sum_A\xi_AD_A
  +\frac12\sum_A D_AD_A
  \\
  &-\frac14\sum_A f_{A\mu\nu}f_A^{\mu\nu}
  -\frac12\sum_A\left(\bar\lambda_A(\sl{D}\lambda)_A\right)
  +\frac{g^2\theta}{64\pi^2}\epsilon_{\mu\nu\rho\sigma}
    \sum_Af_A^{\mu\nu}f_A^{\rho\sigma}\,.
  \end{aligned}
  \tag{27.4.1}\label{eq:27.4.1}
\end{equation}
Here \(f(\phi)\) is the superpotential, a gauge-invariant complex function of
\(\phi_n\) (but not of \(\phi_n^*\)) which the condition of renormalizability
requires to be a cubic polynomial; the \(\xi_A\) are constants which gauge
invariance requires to vanish except where \(t_A\) is a \(U(1)\) generator;
the gauge-covariant derivatives are
\begin{equation}
  D_\mu\psi_L
  \equiv\partial_\mu\psi_L-i\sum_A t_A V_{A\mu}\psi_L\,,
  \tag{27.4.2}\label{eq:27.4.2}
\end{equation}
\begin{equation}
  D_\mu\phi
  \equiv\partial_\mu\phi-i\sum_A t_A V_{A\mu}\phi\,,
  \tag{27.4.3}\label{eq:27.4.3}
\end{equation}
\begin{equation}
  (D_\mu\lambda)_A
  =\partial_\mu\lambda_A+\sum_{BC}C_{ABC}V_{B\mu}\lambda_C\,,
  \tag{27.4.4}\label{eq:27.4.4}
\end{equation}
and \(f_{A\mu\nu}\) is the gauge-covariant gauge field-strength tensor
\begin{equation}
  f_{A\mu\nu}
  =\partial_\mu V_{A\nu}-\partial_\nu V_{A\mu}
  +\sum_{BC}C_{ABC}V_{B\mu}V_{C\nu}\,.
  \tag{27.4.5}\label{eq:27.4.5}
\end{equation}

The auxiliary fields enter quadratically, with field-independent constants as
the coefficients of the second-order term, so they can be eliminated by
setting them equal to the values at which the Lagrangian density is
stationary:
\begin{equation}
  \mathcal{F}_n
  =-\left(\frac{\partial f(\phi)}{\partial\phi_n}\right)^*\,,
  \tag{27.4.6}\label{eq:27.4.6}
\end{equation}
\begin{equation}
  D_A=\xi_A+\sum_{nm}\phi_n^*(t_A)_{nm}\phi_m\,.
  \tag{27.4.7}\label{eq:27.4.7}
\end{equation}
Using these back in Eq.~\eqref{eq:27.4.1}, the Lagrangian density becomes
\begin{equation}
  \begin{aligned}
  \mathcal{L}={}&
  -\sum_n(D_\mu\phi)_n^*(D^\mu\phi)_n
  -\frac12\sum_n
    \left(\bar\psi_{nL}\gamma^\mu(D_\mu\psi_L)_n\right)
  +\frac12\sum_n
    \left(\overline{(D_\mu\psi_L)_n}\gamma^\mu\psi_{nL}\right)
  \\
  &-\frac12\sum_{nm}
    \frac{\partial^2f(\phi)}{\partial\phi_n\partial\phi_m}
    \left(\psi_{nL}^{\mathrm T}\epsilon\psi_{mL}\right)
  -\frac12\sum_{nm}
    \left(
      \frac{\partial^2f(\phi)}{\partial\phi_n\partial\phi_m}
    \right)^*
    \left(\psi_{nL}^{\mathrm T}\epsilon\psi_{mL}\right)^*
  \\
  &-\sum_n\left|
    \frac{\partial f(\phi)}{\partial\phi_n}
  \right|^2
  +i\sqrt{2}\sum_{Anm}
    \left(\bar\psi_{nL}(t_A)_{nm}\lambda_A\right)\phi_m
  \\
  &-i\sqrt{2}\sum_{Anm}
    \phi_n^*\left(\bar\lambda_A(t_A)_{nm}\psi_{mL}\right)
  -\frac12\sum_A
    \left(\xi_A+\sum_{nm}\phi_n^*(t_A)_{nm}\phi_m\right)^2
  \\
  &-\frac14\sum_Af_{A\mu\nu}f_A^{\mu\nu}
  -\frac12\sum_A\left(\bar\lambda_A(\sl{D}\lambda)_A\right)
  +\frac{g^2\theta}{64\pi^2}\epsilon_{\mu\nu\rho\sigma}
    \sum_Af_A^{\mu\nu}f_A^{\rho\sigma}\,.
  \end{aligned}
  \tag{27.4.8}\label{eq:27.4.8}
\end{equation}
Lorentz invariance requires the fields \(\psi_{nL}\), \(\lambda_A\), and
\(f_{A\mu\nu}\) to have vanishing vacuum expectation values, while the
tree-value vacuum expectation values of the \(\phi_n\) are at the minimum of
the potential
\begin{equation}
  V(\phi)
  =\sum_n\left|\frac{\partial f(\phi)}{\partial\phi_n}\right|^2
  +\frac12\sum_A
    \left(\xi_A+\sum_{nm}\phi_n^*(t_A)_{nm}\phi_m\right)^2\,.
  \tag{27.4.9}\label{eq:27.4.9}
\end{equation}
This potential is positive, so if there is a set of field values at which
\(V(\phi)\) vanishes, then this is automatically also a minimum of the
potential. For \(V(\phi)\) to vanish at some field value
\(\phi_n=\phi_{n0}\), it is necessary and sufficient that
\begin{equation}
  \mathcal{F}_{n0}
  =-\left[
    \frac{\partial f(\phi)}{\partial\phi_n}
  \right]^*_{\phi=\phi_0}
  =0
  \tag{27.4.10}\label{eq:27.4.10}
\end{equation}
and
\begin{equation}
  D_{A0}
  =\xi_A+\sum_{nm}\phi_{n0}^*(t_A)_{nm}\phi_{m0}
  =0\,.
  \tag{27.4.11}\label{eq:27.4.11}
\end{equation}
This in turn is the necessary and sufficient condition for supersymmetry not
to be spontaneously broken, since Eq.~\eqref{eq:26.3.15} gives
\(\langle\delta\psi_{nL}\rangle_{\mathrm{VAC}}
=\sqrt{2}\langle\mathcal{F}_n\rangle_{\mathrm{VAC}}\alpha_L\), and
Eq.~\eqref{eq:26.2.16} gives
\(\langle\delta\lambda_A\rangle_{\mathrm{VAC}}
=i\langle D_A\rangle_{\mathrm{VAC}}\gamma_5\alpha\).

It is worth stressing here that spontaneous symmetry breaking is more
difficult for supersymmetry than for other symmetries. For most symmetries of
the action there will be field configurations at which the symmetry is
unbroken and the potential is stationary, but the symmetry will nevertheless
be spontaneously broken if none of these configurations are minima of the
potential. In contrast, any supersymmetric field configuration gives the
potential a value of zero, which is necessarily lower than the value of the
potential for any non-supersymmetric configuration, so the existence of
\textit{any} supersymmetric field configuration insures that supersymmetry is
unbroken. As we will see in Section 27.6, this conclusion goes beyond the tree
approximation used in this section; it is unaffected by corrections of any
finite order in perturbation theory.

It may appear that Eqs.~\eqref{eq:27.4.10} and \eqref{eq:27.4.11} impose too
many conditions on the scalar fields to expect a solution, without some
fine-tuning of the superpotential. However, for a gauge group of
dimensionality \(D\), the superpotential \(f(\phi)\) is subject to the \(D\)
constraints
\begin{equation}
  \sum_m\frac{\partial f(\phi)}{\partial\phi_m}(t_A\phi)_m=0\,,
  \tag{27.4.12}\label{eq:27.4.12}
\end{equation}
for all \(A\) and all \(\phi\). Hence if \(\phi\) has \(N\) independent
components, then the number of \textit{independent} conditions
\eqref{eq:27.4.10} is \(N-D\), while the number of conditions
\eqref{eq:27.4.11} is \(D\), so there are just \(N\) conditions altogether.
With the number of conditions equal to the number of free variables, it is
likely to find solutions for generic superpotentials. In fact, it is more
usual to find solutions than not. For instance, for chiral scalar superfields
in a non-trivial representation of a semi-simple gauge group we have
\(\xi_A=0\), while \(f(\phi)\) can have no terms linear in the \(\phi_n\),
so both Eqs.~\eqref{eq:27.4.10} and \eqref{eq:27.4.11} are satisfied for
\(\phi_{n0}=0\). There may be other solutions of Eqs.~\eqref{eq:27.4.10} and
\eqref{eq:27.4.11} which break the gauge symmetries, but in such a theory
supersymmetry cannot be broken, at least not in the tree approximation, and,
as we shall see in Section 27.6, not in any order of perturbation theory.

More generally, it is easy to see that, even if the gauge group has \(U(1)\)
factors and even if the superpotential involves gauge-invariant superfields,
if there exists a set of scalar field values \(\phi_{n0}\) that satisfy
Eq.~\eqref{eq:27.4.10}, then there is another that satisfies both
Eq.~\eqref{eq:27.4.10} and Eq.~\eqref{eq:27.4.11}, provided only that the
Fayet--Iliopoulos constants \(\xi_A\) all vanish. To show this, we note that
since the superpotential \(f(\phi)\) does not involve \(\phi^*\), it is
invariant not only under ordinary gauge transformations
\(\phi\longrightarrow\exp(i\sum_A\Lambda_A t_A)\phi\), with
\(\Lambda_A\) arbitrary real numbers, but also under transformations with
\(\Lambda_A\) arbitrary complex numbers. Under all these transformations the
\(\mathcal{F}\)-terms in Eq.~\eqref{eq:27.4.10} transform linearly, so if
\(\phi_0\) satisfies Eq.~\eqref{eq:27.4.10}, then so does
\(\phi^\Lambda\equiv\exp(i\sum_A\Lambda_A t_A)\phi_0\). On the other hand,
the scalar product \([\phi^\dagger\phi]\) is not invariant under
transformations with \(\Lambda_A\) complex, but
\([\phi^{\Lambda\dagger}\phi^\Lambda]\) does remain real and positive for
complex \(\Lambda_A\), so it is bounded below, and therefore has a minimum.
For \(\xi_A=0\), the condition that
\([\phi^{\Lambda\dagger}\phi^\Lambda]\) be stationary at this minimum is just
that \(\phi^\Lambda\) should satisfy Eq.~\eqref{eq:27.4.11}. We see then that
in the absence of Fayet--Iliopoulos \(D\)-terms, the question of whether
supersymmetry is unbroken in gauge theories is entirely a matter of whether
the superpotential allows solutions of Eq.~\eqref{eq:27.4.10}. The same
result applies even in non-renormalizable
theories~\hyperref[ch27-ref-3]{[3]}.

Now let us assume that there is a set of values \(\phi_{n0}\) at which
\(V(\phi_0)=0\), so that supersymmetry is unbroken. The spin 0 degrees of
freedom are described by a shifted field
\begin{equation}
  \varphi_n=\phi_n-\phi_{n0}\,.
  \tag{27.4.13}\label{eq:27.4.13}
\end{equation}
There is then a cross-term between \(\varphi\) and the gauge fields, arising
from the first term in Eq.~\eqref{eq:27.4.1}:
\[
  2\sum_{nA}\operatorname{Im}
  \left(\partial_\mu\varphi_n(t_A\phi_0)_n^*\right)V_A^\mu\,.
\]
As shown in Section 21.1, it is always possible to eliminate this term by
adopting a `unitarity gauge,' in which \(\phi_n\) satisfies a constraint that
makes this term vanish:
\begin{equation}
  \sum_n\operatorname{Im}
  \left(\phi_n(t_A\phi_0)_n^*\right)=0\,.
  \tag{27.4.14}\label{eq:27.4.14}
\end{equation}
This will have the effect of eliminating the Goldstone bosons associated with
broken gauge symmetries.

We shall now work out the masses of the particles of spin 0, spin \(1/2\),
and spin 1 that arise in this theory if supersymmetry is unbroken, taking into
account the possibility that the gauge symmetries may be spontaneously
broken.

\subsubsection*{Spin 0}

Because \(\partial f(\phi)/\partial\phi_n\) and
\(\xi_A+\sum_{nm}\phi_n^*(t_A)_{nm}\phi_m\) must both vanish at
\(\phi_n=\phi_{n0}\), the terms in \(V(\phi)\) of second order in
\(\varphi_n\equiv\phi_n-\phi_{n0}\) and/or \(\varphi_n^*\) are of the form
\begin{equation}
  \begin{aligned}
  V_{\mathrm{quad}}(\phi)={}&
  \sum_{nm}(\mathcal{M}^*\mathcal{M})_{nm}\varphi_n^*\varphi_m
  +\sum_{Anm}(t_A\phi_0)_n(t_A\phi_0)_m^*
    \varphi_n^*\varphi_m
  \\
  &+\frac12\sum_{Anm}
    (t_A\phi_0)_n^*(t_A\phi_0)_m^*\varphi_n\varphi_m
  +\frac12\sum_{Anm}
    (t_A\phi_0)_n(t_A\phi_0)_m\varphi_n^*\varphi_m^*\,,
  \end{aligned}
  \tag{27.4.15}\label{eq:27.4.15}
\end{equation}
where \(\mathcal{M}\) is the complex symmetric matrix \eqref{eq:26.4.11}:
\[
  \mathcal{M}_{nm}
  \equiv
  \left(
    \frac{\partial^2f(\phi)}{\partial\phi_n\partial\phi_m}
  \right)_{\phi=\phi_0}.
\]
This can be written as
\begin{equation}
  V_{\mathrm{quad}}
  =\frac12
  \begin{bmatrix}\varphi\\ \varphi^*\end{bmatrix}^{\!\dagger}
  M_0^2
  \begin{bmatrix}\varphi\\ \varphi^*\end{bmatrix},
  \tag{27.4.16}\label{eq:27.4.16}
\end{equation}
where \(M_0^2\) is the block matrix
\begin{equation}
  M_0^2=
  \begin{bmatrix}
  \mathcal{M}^*\mathcal{M}
    +\sum_A(t_A\phi_0)(t_A\phi_0)^\dagger
  &
  \sum_A(t_A\phi_0)(t_A\phi_0)^{\mathrm T}
  \\
  \sum_A(t_A\phi_0)^*(t_A\phi_0)^\dagger
  &
  \mathcal{M}\mathcal{M}^*
    +\sum_A(t_A\phi_0)^*(t_A\phi_0)^{\mathrm T}
  \end{bmatrix}.
  \tag{27.4.17}\label{eq:27.4.17}
\end{equation}
Now we must look for the eigenvalues of this mass-squared matrix.
Differentiating Eq.~\eqref{eq:27.4.12} with respect to \(\phi_n\) gives
\begin{equation}
  \sum_m
  \frac{\partial^2f(\phi)}{\partial\phi_n\partial\phi_m}
  (t_A\phi)_m
  +\sum_m\frac{\partial f(\phi)}{\partial\phi_m}(t_A)_{mn}
  =0\,.
  \tag{27.4.18}\label{eq:27.4.18}
\end{equation}
But as we have seen, \(\partial f(\phi)/\partial\phi_m\) vanishes at
\(\phi=\phi_0\), so by setting \(\phi\) at this value in
Eq.~\eqref{eq:27.4.18}, we find
\begin{equation}
  \sum_m\mathcal{M}_{nm}(t_A\phi_0)_m=0\,.
  \tag{27.4.19}\label{eq:27.4.19}
\end{equation}
It follows that
\[
  M_0^2
  \begin{bmatrix}
    t_B\phi_0\\
    \mathord{\pm}(t_B\phi_0)^*
  \end{bmatrix}
  =
  \sum_A
  \left(\phi_0^\dagger[t_A t_B\mathord{\pm}t_B t_A]\phi_0\right)
  \begin{bmatrix}
    t_B\phi_0\\
    \mathord{\pm}(t_B\phi_0)^*
  \end{bmatrix}.
\]
But the vanishing of \(D_A\) at \(\phi=\phi_0\) and the global gauge
invariance of \(\xi_A\) tell us that
\begin{equation}
  \left(\phi_0^\dagger[t_A,t_B]\phi_0\right)
  =i\sum_C C_{ABC}\left(\phi_0^\dagger t_C\phi_0\right)
  =-i\left(\phi_0^\dagger\phi_0\right)
    \sum_C C_{ABC}\xi_C
  =0\,.
  \tag{27.4.20}\label{eq:27.4.20}
\end{equation}
The matrix \eqref{eq:27.4.17} therefore has a pair of eigenvectors for each
gauge symmetry
\begin{equation}
  u=
  \begin{bmatrix}
    \sum_Bc_Bt_B\phi_0\\
    \sum_Bc_B(t_B\phi_0)^*
  \end{bmatrix},
  \qquad
  v=
  \begin{bmatrix}
    \sum_Bc_Bt_B\phi_0\\
    -\sum_Bc_B(t_B\phi_0)^*
  \end{bmatrix},
  \tag{27.4.21}\label{eq:27.4.21}
\end{equation}
for which
\begin{equation}
  M_0^2u=\mu^2u\,,
  \qquad
  M_0^2v=0\,,
  \tag{27.4.22}\label{eq:27.4.22}
\end{equation}
where \(\mu^2\) and \(c_A\) are any real solutions of the eigenvalue
problem\footnote{The presence of a factor \(1/2\) in Eq.~\eqref{eq:21.1.17}
which does not appear in Eq.~\eqref{eq:27.4.23} is due to a difference in
the way that the scalar fields are normalized.}
\begin{equation}
  \sum_B
  \left(\phi_0^\dagger\{t_A,t_B\}\phi_0\right)c_B
  =\mu^2c_A\,,
  \tag{27.4.23}\label{eq:27.4.23}
\end{equation}
with the exception that if the eigenvalue \(\mu^2\) vanishes then
\(\sum_Bc_Bt_B\phi_0=0\), so that the eigenvectors \(u\) and \(v\) are
absent. The massless particles associated with the \(v\) eigenvectors are
Goldstone bosons, which are eliminated from the physical spectrum by the
unitarity gauge condition \eqref{eq:27.4.14}. In addition to these mass
eigenstates, there is another set that are orthogonal to all the \(u\)s and
\(v\)s and that therefore take the form
\begin{equation}
  w_\pm=
  \begin{bmatrix}
    \zeta\\
    \mathord{\pm}\zeta^*
  \end{bmatrix},
  \tag{27.4.24}\label{eq:27.4.24}
\end{equation}
where
\begin{equation}
  \sum_n(t_A\phi_0)_n^*\zeta_n=0\,.
  \tag{27.4.25}\label{eq:27.4.25}
\end{equation}
Eq.~\eqref{eq:27.4.19} shows that the space of \(\zeta\)s satisfying
Eq.~\eqref{eq:27.4.25} is invariant under multiplication with the Hermitian
matrix \(\mathcal{M}^\dagger\mathcal{M}\), so it is spanned by the
eigenvectors of this matrix, satisfying
\begin{equation}
  \mathcal{M}^\dagger\mathcal{M}\zeta=m^2\zeta
  \tag{27.4.26}\label{eq:27.4.26}
\end{equation}
with \(m^2\) a set of real positive (or zero) eigenvalues.
Eq.~\eqref{eq:27.4.26} and its complex conjugate together with
Eq.~\eqref{eq:27.4.25} show that \(w_\pm\) are eigenvectors of \(M_0^2\)
with eigenvalues \(m^2\):
\begin{equation}
  M_0^2w_\pm=m^2w_\pm\,.
  \tag{27.4.27}\label{eq:27.4.27}
\end{equation}
We thus have \textit{two} self-charge-conjugate spinless bosons of each mass
\(m\) satisfying Eq.~\eqref{eq:27.4.27}, and one self-charge-conjugate
spinless boson of each non-zero mass \(\mu\) satisfying
Eq.~\eqref{eq:27.4.23}.

\subsubsection*{Spin \(1/2\)}

The fermion masses arise from the non-derivative terms in
Eq.~\eqref{eq:27.4.8} that are of second order in the fermion fields
\(\psi_n\) and \(\lambda_A\):
\begin{equation}
  \mathcal{L}_{1/2}
  =-\frac12\sum_{nm}\mathcal{M}_{nm}
    \left(\psi_{nL}^{\mathrm T}\epsilon\psi_{mL}\right)
  -i\sqrt{2}\sum_{Am}(t_A\phi_0)_m^*
    \left(\lambda_{LA}^{\mathrm T}\epsilon\psi_{mL}\right)
  +\mathrm{H.c.}
  \tag{27.4.28}\label{eq:27.4.28}
\end{equation}
We saw in Section 26.4 that if the fermion mass term in the Lagrangian for a
column \(\chi\) of Majorana spinor fields is put in the form
\begin{equation}
  \mathcal{L}_{1/2}
  =-\frac12\left(\chi_L^{\mathrm T}\epsilon M\chi_L\right)
  +\mathrm{H.c.}\,,
  \tag{27.4.29}\label{eq:27.4.29}
\end{equation}
then the fermion squared masses are the eigenvalues of the Hermitian matrix
\(M^\dagger M\). Eq.~\eqref{eq:27.4.28} gives the elements of the matrix
\(M\) here as
\begin{equation}
  M_{nm}=\mathcal{M}_{nm}\,,
  \qquad
  M_{nA}=M_{An}=i\sqrt{2}(t_A\phi_0)_n^*\,,
  \qquad
  M_{AB}=0\,,
  \tag{27.4.30}\label{eq:27.4.30}
\end{equation}
for which, using Eqs.~\eqref{eq:27.4.19} and \eqref{eq:27.4.20},
\begin{equation}
  \begin{aligned}
  (M^\dagger M)_{nm}
  &=(\mathcal{M}^\dagger\mathcal{M})_{nm}
    +2\sum_A(t_A\phi_0)_n(t_A\phi_0)_m^*\,,
  \\
  (M^\dagger M)_{nA}
  &=(M^\dagger M)_{An}=0\,,
  \\
  (M^\dagger M)_{AB}
  &=2\left(\phi_0^\dagger t_Bt_A\phi_0\right)
   =\left(\phi_0^\dagger\{t_B,t_A\}\phi_0\right).
  \end{aligned}
  \tag{27.4.31}\label{eq:27.4.31}
\end{equation}
The eigenvectors of the matrix \eqref{eq:27.4.30} are of three types. First
are those of the form
\begin{equation}
  z=
  \begin{bmatrix}
    \zeta\\
    0
  \end{bmatrix},
  \tag{27.4.32}\label{eq:27.4.32}
\end{equation}
with eigenvalues \(m^2\), where \(\zeta_n\) and \(m^2\) are any eigenvectors
and corresponding eigenvalues of \(\mathcal{M}^\dagger\mathcal{M}\). Next
are those of the form
\begin{equation}
  g=
  \begin{bmatrix}
    0\\
    c
  \end{bmatrix},
  \tag{27.4.33}\label{eq:27.4.33}
\end{equation}
with eigenvalues \(\mu^2\), where \(c_B\) and \(\mu^2\) are any eigenvectors
and eigenvalues of the matrix
\(\left(\phi_0^\dagger\{t_B,t_A\}\phi_0\right)\). Finally, there are those
of the form
\begin{equation}
  h=
  \begin{bmatrix}
    \sum_Bc_Bt_B\phi_0\\
    0
  \end{bmatrix},
  \tag{27.4.34}\label{eq:27.4.34}
\end{equation}
with eigenvalues \(\mu^2\), where \(c_B\) and \(\mu^2\) are again any
eigenvectors and eigenvalues of the matrix
\(\left(\phi_0^\dagger\{t_B,t_A\}\phi_0\right)\). The only exception is
that the eigenvectors \(c\) of this matrix with eigenvalue zero have
\(\sum_Ac_At_A\phi_0=0\), corresponding to unbroken symmetries, so that in
this case the vector \eqref{eq:27.4.34} vanishes and we have only the
eigenvector \eqref{eq:27.4.33}. Thus there is one Majorana fermion of each
mass \(m\) satisfying Eq.~\eqref{eq:27.4.26}, two Majorana fermions of each
non-zero mass \(\mu\) satisfying Eq.~\eqref{eq:27.4.22}, and one Majorana
fermion of zero mass for each unbroken gauge symmetry.

\subsubsection*{Spin 1}

The mass terms in the Lagrangian for the gauge fields arise from the part of
the first term in Eq.~\eqref{eq:27.4.1} that is of second order in the gauge
field \(V_A^\mu\):
\begin{equation}
  \mathcal{L}_V
  =-\sum_{nAB}(t_A\phi_0)_n^*(t_B\phi_0)_nV_{A\mu}V_B^\mu\,.
  \tag{27.4.35}\label{eq:27.4.35}
\end{equation}
Since the fields \(V_{A\mu}\) are real, their mass-squared matrix is the
matrix in Eq.~\eqref{eq:27.4.23}:
\begin{equation}
  (\mu^2)_{AB}
  =\left(\phi_0^\dagger\{t_B,t_A\}\phi_0\right).
  \tag{27.4.36}\label{eq:27.4.36}
\end{equation}
There is one spin 1 particle of mass \(\mu\) for each eigenvalue \(\mu^2\) of
the matrix \eqref{eq:27.4.36}.

Putting this all together, we see that for each eigenvalue \(m^2\) of the
matrix \(\mathcal{M}^*\mathcal{M}\) there are two self-charge-conjugate
spinless particles and one Majorana fermion of mass \(m\); for each non-zero
eigenvalue of the matrix \(\mu^2_{AB}\) there are one self-charge-conjugate
spinless boson, two Majorana fermions, and one self-charge-conjugate spin 1
boson of mass \(\mu\); and for each zero eigenvalue of this matrix there is
one Majorana fermion and one self-charge-conjugate spin 1 boson of zero mass.
It is not surprising that the particle multiplets for each zero or non-zero
mass are just the same as we found by direct use of the supersymmetry algebra
in Sections 25.4 and 25.5. What is a bit surprising is that the masses of the
gauge and chiral particles are unaffected by each other. The masses \(m\)
that are given by eigenvalues of \((\mathcal{M}^*\mathcal{M})_{nm}\) and the
particles with these masses are just what they would be in a theory of chiral
superfields without gauge superfields, and the masses \(\mu\) that are given
by eigenvalues of \(\mu^2_{AB}\) and the particles with these masses are just
what they would be in a theory of gauge superfields with no chiral
superfields.

For future use in Section 27.9 we will now apply the method described in
Section 26.7 to construct the supersymmetry current for the supersymmetric
gauge Lagrangian \eqref{eq:27.4.1}. In the gauge used earlier, an
infinitesimal supersymmetry transformation changes \(V_A\), \(\lambda_A\),
and \(D_A\) by the amounts \eqref{eq:27.3.4}--\eqref{eq:27.3.6}. Adding the
Noether supersymmetry current given by Eq.~\eqref{eq:26.7.2} for these fields
to the Noether current for \(\phi_n\), \(\psi_n\), and \(\mathcal{F}_n\)
already given in Eq.~\eqref{eq:26.7.8}, with derivatives replaced by
gauge-invariant derivatives, gives the total Noether supersymmetry current:
\begin{equation}
  \begin{aligned}
  N^\mu={}&
  \sum_Af_A^{\mu\nu}\gamma_\nu\lambda_A
  -\frac18\sum_Af_{A\rho\sigma}
    [\gamma^\rho,\gamma^\sigma]\gamma^\mu\lambda_A
  -\frac12i\sum_AD_A\gamma_5\gamma^\mu\lambda_A
  \\
  &+\frac1{\sqrt{2}}\sum_n\left[
    2(D^\mu\phi)_n^*\psi_{nL}
    +2(D^\mu\phi)_n\psi_{nR}
    +(\sl{D}\phi)_n\gamma^\mu\psi_{nR}
  \right.
  \\
  &\hspace{36mm}\left.
    +(\sl{D}\phi)_n^*\gamma^\mu\psi_{nL}
    -\mathcal{F}_n\gamma^\mu\psi_{nR}
    -\mathcal{F}_n^*\gamma^\mu\psi_{nL}
  \right].
  \end{aligned}
  \tag{27.4.37}\label{eq:27.4.37}
\end{equation}
This is not the supersymmetry current, because the Lagrangian density is not
invariant under supersymmetry; instead, its change is the derivative
\begin{equation}
  \delta\mathcal{L}=\partial_\mu\left(\bar\alpha K^\mu\right),
  \tag{27.4.38}\label{eq:27.4.38}
\end{equation}
where\footnote{The easiest way to calculate the change in the term
\(\left[\Phi^\dagger\exp(-2\sum_A t_AV_A)\Phi\right]_D\) is to calculate the
\(\lambda\)-component of
\(\Phi^\dagger\exp(-2\sum_A t_AV_A)\Phi\) and use
Eq.~\eqref{eq:26.2.17}. In this way of doing the calculation, the important
term on the second line of the right-hand side of Eq.~\eqref{eq:27.4.39}
arises from the \(\lambda\)-component of
\(\exp(-2\sum_A t_AV_A)\).}
\begin{equation}
  \begin{aligned}
  K^\mu={}&
  \frac12i\sum_A
    \epsilon^{\rho\sigma\mu\nu}f_{A\rho\sigma}
    \gamma_\nu\gamma_5\lambda_A
  +\frac18\sum_A[\gamma^\rho,\gamma^\sigma]
    \gamma^\mu\lambda_Af_{A\rho\sigma}
  +\frac12i\sum_AD_A\gamma_5\gamma^\mu\lambda_A
  \\
  &-i\sum_{Anm}(t_A)_{nm}
    \gamma_5\gamma^\mu\lambda_A\phi_n^*\phi_m
  \\
  &+\frac1{\sqrt{2}}\sum_n\gamma^\mu\left[
    -(\sl{D}\phi)_n\psi_{nR}
    -(\sl{D}\phi)_n^*\psi_{nL}
    +\mathcal{F}_n^*\psi_{nL}
    +\mathcal{F}_n\psi_{nR}
  \right.
  \\
  &\hspace{39mm}\left.
    +2\left(\frac{\partial f(\phi)}{\partial\phi_n}\right)^*\psi_{nL}
    +2\left(\frac{\partial f(\phi)}{\partial\phi_n}\right)\psi_{nR}
  \right].
  \end{aligned}
  \tag{27.4.39}\label{eq:27.4.39}
\end{equation}
The first two terms are derived using the identity \eqref{eq:27.2.5}. Using
the same identity again with Eq.~\eqref{eq:26.7.4} gives the total
supersymmetry current:
\begin{equation}
  \begin{aligned}
  S^\mu=N^\mu+K^\mu
  ={}&
  -\frac14\sum_Af_{A\rho\sigma}
    [\gamma^\rho,\gamma^\sigma]\gamma^\mu\lambda_A
  -i\sum_{Anm}(t_A)_{nm}
    \gamma_5\gamma^\mu\lambda_A\phi_n^*\phi_m
  \\
  &+\frac1{\sqrt{2}}\sum_n\left[
    (\sl{D}\phi)_n\gamma^\mu\psi_{nR}
    +(\sl{D}\phi^*)_n\gamma^\mu\psi_{nL}
  \right.
  \\
  &\hspace{31mm}\left.
    +2\left(\frac{\partial f(\phi)}{\partial\phi_n}\right)
      \gamma^\mu\psi_{nL}
    +2\left(\frac{\partial f(\phi)}{\partial\phi_n}\right)^*
      \gamma^\mu\psi_{nR}
  \right].
  \end{aligned}
  \tag{27.4.40}\label{eq:27.4.40}
\end{equation}

\begin{center}
\(\ast\qquad\ast\qquad\ast\)
\end{center}

In Section 26.8 we considered a class of supersymmetric theories with a
superpotential \(f(\Phi)\) that has an arbitrary dependence on a set of
left-chiral scalar superfields \(\Phi_n\) but not their derivatives, and with
a Kahler potential \(K(\Phi,\Phi^*)\) that has an arbitrary dependence on the
\(\Phi_n\) and \(\Phi_n^*\), but not their derivatives. We can extend the
same considerations to gauge theories, again with the dependence of the
Lagrangian on the chiral superfields limited only by supersymmetry, but
without introducing new superderivatives or spacetime derivatives. The
renormalizable Lagrangian density is then replaced with
\begin{equation}
  \begin{aligned}
  \mathcal{L}={}&
  \frac12\left[
    K\left(
      \Phi,\Phi^\dagger
      \exp\left(-2\sum_A t_AV_A\right)
    \right)
  \right]_D
  +2\operatorname{Re}[f(\Phi)]_{\mathcal{F}}
  \\
  &-\frac12\operatorname{Re}\sum_{AB}
  \left[
    h_{AB}(\Phi)
    \left(W_{AL}^{\mathrm T}\epsilon W_{BL}\right)
  \right]_{\mathcal{F}}\,,
  \end{aligned}
  \tag{27.4.41}\label{eq:27.4.41}
\end{equation}
where \(h_{AB}(\Phi)\) is a new function of the \(\Phi_n\), but not of the
\(\Phi_n^*\) or derivatives.

The chiral gauge and scalar superfields are given by the expansions
\eqref{eq:26.3.21} and \eqref{eq:27.3.15}:
\[
  \begin{aligned}
  W_{AL}(x,\theta)={}&
  \lambda_{AL}(x_+)
  +\frac12\gamma^\mu\gamma^\nu\theta_Lf_{A\mu\nu}(x_+)
  +(\theta_L^{\mathrm T}\epsilon\theta_L)
    \sl{\partial}\lambda_{AR}(x_+)
  -i\theta_LD_A(x_+),
  \\
  \Phi_n(x,\theta)={}&
  \phi_n(x_+)
  -\sqrt{2}\left(\theta_L^{\mathrm T}\epsilon\psi_{nL}(x_+)\right)
  +\mathcal{F}_n(x_+)
    \left(\theta_L^{\mathrm T}\epsilon\theta_L\right),
  \end{aligned}
\]
where \(x_+^\mu\) is the shifted coordinate \eqref{eq:26.3.23}. The terms of
second order in \(\theta_L\) (and independent of \(\theta_R\)) in
\(\sum_{AB}h_{AB}(\Phi)(W_{AL}^{\mathrm T}\epsilon W_{BL})\) are then
\[
  \begin{aligned}
  -\left[
    \sum_{AB}h_{AB}(\Phi)
    \left(W_{AL}^{\mathrm T}\epsilon W_{BL}\right)
  \right]_{\theta_L^2}
  ={}&
  (\theta_L^{\mathrm T}\epsilon\theta_L)
  \sum_{AB}\left(\lambda_{AL}^{\mathrm T}\epsilon\lambda_{BL}\right)
  \left[
    \frac12\sum_{nm}
      \left(\psi_{nL}^{\mathrm T}\epsilon\psi_{mL}\right)
      \frac{\partial^2h_{AB}(\phi)}
           {\partial\phi_n\partial\phi_m}
  \right.
  \\
  &\hspace{49mm}\left.
    -\sum_n\mathcal{F}_n
      \frac{\partial h_{AB}(\phi)}{\partial\phi_n}
  \right]
  \\
  &+(\theta_L^{\mathrm T}\epsilon\theta_L)
  \sum_{AB}h_{AB}(\phi)
  \left[
    -\left(\bar\lambda_A\sl{\partial}
      (1-\gamma_5)\lambda_B\right)
    -\frac12f_{A\mu\nu}f_B^{\mu\nu}
  \right.
  \\
  &\hspace{47mm}\left.
    +\frac{i}{4}\epsilon_{\mu\nu\rho\sigma}
      f_A^{\mu\nu}f_B^{\rho\sigma}
    +D_AD_B
  \right]
  \\
  &+\sqrt{2}\sum_{ABn}
  \frac{\partial h_{AB}(\phi)}{\partial\phi_n}
  \left(\theta_L^{\mathrm T}\epsilon\psi_{nL}\right)
  \left[
    -\left(\lambda_{BL}^{\mathrm T}
      \epsilon\gamma^\mu\gamma^\nu\theta_L\right)f_{A\mu\nu}
  \right.
  \\
  &\hspace{54mm}\left.
    +2i\left(\lambda_{BL}^{\mathrm T}\epsilon\theta_L\right)
  \right],
  \end{aligned}
\]
with all fields now understood to be evaluated at \(x^\mu\) rather than
\(x_+^\mu\). (The first and second terms on the right-hand side are taken
from Eqs.~\eqref{eq:26.4.4} and \eqref{eq:27.3.16}, respectively.) Also, by
writing \(\theta_{L\alpha}\theta_{L\beta}\) as
\(\tfrac12\epsilon_{\alpha\beta}
(\theta_L^{\mathrm T}\epsilon\theta_L)\), the third term on the right-hand
side can also be expressed as proportional to
\((\theta_L^{\mathrm T}\epsilon\theta_L)\):
\[
  \begin{aligned}
  &\left(\theta_L^{\mathrm T}\epsilon\psi_{nL}\right)
  \left[
    \left(\bar\psi_B\gamma^\mu\gamma^\nu\theta_L\right)f_{A\mu\nu}
    -2i\left(\bar\psi_B\theta_L\right)
  \right]
  \\
  &\qquad
  =\frac12\left(\theta_L^{\mathrm T}\epsilon\theta_L\right)
  \left[
    \left(\bar\psi_B\gamma^\mu\gamma^\nu\psi_{nL}\right)
    -2i\left(\bar\psi_B\psi_{nL}\right)D_A
  \right].
  \end{aligned}
\]
The \(\mathcal{F}\)-term is the coefficient of
\((\theta_L^{\mathrm T}\epsilon\theta_L)\), so
\[
  \begin{aligned}
  -\left[
    \sum_{AB}h_{AB}(\Phi)
    \left(W_{AL}^{\mathrm T}\epsilon W_{BL}\right)
  \right]_{\mathcal{F}}
  ={}&
  \sum_{AB}\left(\lambda_{AL}^{\mathrm T}\epsilon\lambda_{BL}\right)
  \\
  &\quad{}\times
  \left[
    \frac12\sum_{nm}
      \left(\psi_{nL}^{\mathrm T}\epsilon\psi_{mL}\right)
      \frac{\partial^2h_{AB}(\phi)}
           {\partial\phi_n\partial\phi_m}
  \right.
  \\
  &\hspace{48mm}\left.
    -\sum_n\mathcal{F}_n
      \frac{\partial h_{AB}(\phi)}{\partial\phi_n}
  \right]
  \\
  &+\sum_{AB}h_{AB}(\phi)
  \left[
    -\left(\bar\lambda_A\sl{\partial}
      (1-\gamma_5)\lambda_B\right)
    -\frac12f_{A\mu\nu}f_B^{\mu\nu}
  \right.
  \\
  &\hspace{47mm}\left.
    +\frac{i}{4}\epsilon_{\mu\nu\rho\sigma}
      f_A^{\mu\nu}f_B^{\rho\sigma}
    +D_AD_B
  \right]
  \\
  &+\frac{\sqrt{2}}2\sum_{ABn}
  \frac{\partial h_{AB}(\phi)}{\partial\phi_n}
  \\
  &\quad{}\times
  \left[
    -\left(\bar\psi_B\gamma^\mu\gamma^\nu\psi_{nL}\right)f_{A\mu\nu}
    +2i\left(\bar\psi_B\psi_{nL}\right)D_A
  \right].
  \end{aligned}
\]
The other terms in Eq.~\eqref{eq:27.4.41} are just given by the
gauge-invariant version of the Lagrangian density \eqref{eq:26.8.6}. Putting
this together gives the Lagrangian density
\begin{equation}
  \begin{aligned}
  \mathcal{L}={}&
  \operatorname{Re}\sum_{nm}\mathcal{G}_{nm}(\phi,\phi^*)
  \left[
    -\frac12\left(\bar\psi_m\sl{D}(1+\gamma_5)\psi_n\right)
    +\mathcal{F}_n\mathcal{F}_m^*
    -D_\mu\phi_nD^\mu\phi_m^*
  \right]
  \\
  &-2\operatorname{Re}\sum_i
    \frac{\partial K(\phi,\phi^*)}{\partial\phi_i^*}
    D_A(\phi^*t_A)_i
  \\
  &+i\sqrt{2}
    \frac{\partial^2K(\phi,\phi^*)}
         {\partial\phi_i\partial\phi_j^*}
    \left[
      (t_A\phi)_i\bar\psi_j\lambda_{AR}
      -(\phi^*t_A)_j\bar\psi_i\lambda_{AL}
    \right]
  \\
  &-\operatorname{Re}\sum_{nm\ell}
    \frac{\partial^3K(\phi,\phi^*)}
         {\partial\phi_n\partial\phi_m\partial\phi_\ell^*}
    \left(\bar\psi_n\psi_{mL}\right)\mathcal{F}_\ell^*
  \\
  &+\operatorname{Re}\sum_{nm\ell}
    \frac{\partial^3K(\phi,\phi^*)}
         {\partial\phi_n\partial\phi_m\partial\phi_\ell^*}
    \left(\bar\psi_m\gamma^\mu\psi_{\ell R}\right)D_\mu\phi_n
  \\
  &+\frac14\sum_{nm\ell k}
    \frac{\partial^4K(\phi,\phi^*)}
         {\partial\phi_n\partial\phi_m
          \partial\phi_\ell^*\partial\phi_k^*}
    \left(\bar\psi_n\psi_{mL}\right)
    \left(\bar\psi_k\psi_{\ell R}\right)
  \\
  &-\operatorname{Re}\sum_{nm}
    \frac{\partial^2f(\phi)}{\partial\phi_n\partial\phi_m}
    \left(\bar\psi_n\psi_{mL}\right)
  +2\operatorname{Re}\sum_n
    \mathcal{F}_n\frac{\partial f(\phi)}{\partial\phi_n}
  \\
  &+\frac14\operatorname{Re}\sum_{ABnm}
    \left(\bar\lambda_A\lambda_{BL}\right)
    \left(\bar\psi_n\psi_{mL}\right)
    \frac{\partial^2h_{AB}(\phi)}
         {\partial\phi_n\partial\phi_m}
  \\
  &-\frac12\operatorname{Re}\sum_{ABn}
    \left(\bar\lambda_A\lambda_{BL}\right)\mathcal{F}_n
    \frac{\partial h_{AB}(\phi)}{\partial\phi_n}
  \\
  &+\operatorname{Re}\sum_{AB}h_{AB}(\phi)
  \left[
    -\left(\bar\lambda_A\sl{D}\lambda_{BR}\right)
    -\frac14f_{A\mu\nu}f_B^{\mu\nu}
    +\frac{i}{8}\epsilon_{\mu\nu\rho\sigma}
      f_A^{\mu\nu}f_B^{\rho\sigma}
    +\frac12D_AD_B
  \right]
  \\
  &+\frac{\sqrt{2}}4\operatorname{Re}\sum_{ABn}
    \frac{\partial h_{AB}(\phi)}{\partial\phi_n}
    \left[
      -\left(\bar\lambda_B\gamma^\mu\gamma^\nu\psi_{nL}\right)
       f_{A\mu\nu}
      +2i\left(\bar\lambda_B\psi_{nL}\right)D_A
    \right].
  \end{aligned}
  \tag{27.4.42}\label{eq:27.4.42}
\end{equation}
One interesting feature of this result is the appearance of a gaugino mass
in theories with \(\phi_n\)-dependent functions \(h_{AB}(\phi)\), when
supersymmetry is broken by a non-vanishing value of \(\mathcal{F}_n\). This
mechanism is used to generate gaugino masses in some theories of
gravitationally mediated supersymmetry breaking, discussed in Section 31.7.
